% ============================================================================
% Fichier  : partie2/P02_C05_quantifier_fair.tex
% Titre    : Quantifier le Risque Cyber : FAIR
% Auteur   : MINKA MI NGUIDJOÏ Thierry Emmanuel
% Version  : 1.0 -- novembre 2026
% Corrigés : corriges/Corr_C05.tex
% ============================================================================
% Sources : Open FAIR (The Open Group), O-RT 3.0.1 et O-RA 2.0.1 (versions
% vérifiées sur opengroup.org le 05/11/2026) ; taxonomie présentée d'après
% Freund et Jones (2015). Le renommage « vulnérabilité » → « susceptibilité »
% dans O-RT 3 est signalé à l'auteur comme À VÉRIFIER. Estimations calibrées
% : Hubbard et Seiersen (2016).
% Faits du cas (bible) : R-AG-01, fréquence 0,2 à 0,6, ampleur 150 à 900 M
% FCFA ; incident réel 410 M ; plan A1 à A8 = 193 M ; assurance (prime
% 35 M, franchise 100 M, plafond 1 Md).
% Faits nouveaux (reportés dans la bible) : valeurs les plus probables
% 0,35 et 410 M ; ventilation des 410 M (productivité 190, réponse 140,
% remplacement 80) ; après traitement : fréquence PERT(0,05 ; 0,10 ; 0,20),
% ampleur PERT(50 ; 150 ; 400) ; coût récurrent du plan 30 M par an.
% CHIFFRES CALCULÉS par simulation (scratchpad fair_sanaga.py, numpy,
% graine 2027, 100 000 itérations, fréquence PERT tirée à chaque itération,
% nombre d'événements de Poisson, ampleur PERT par événement) :
%   actuel : moyenne 164,6 ; P(L>0) 0,305 ; P50 0 ; P90 579,3 ; P95 733,6 ;
%   P99 1 169,8 ; P(L>500) 0,143 ; P(lambda>1/3) 0,640 ; P(k>=2) 0,054 ;
%   dépassements 100:0,305 250:0,288 500:0,143 750:0,047 1000:0,021.
%   après : moyenne 18,6 ; P(L>0) 0,101 ; P90 60,0 ; P95 174,4 ; P99 283,2 ;
%   P(L>500) 0,0006 ; dépassements 100:0,089 250:0,018 500:0,0006.
%   assurance : remboursement moyen 126,6 (actuel), 8,2 (après).
%   sensibilité, max 1 500 : moyenne 200,9 ; P99 1 507,0 ; P(L>500) 0,184.
%   table des 10 années : graine 7.
% ============================================================================

\chapter{Quantifier le Risque Cyber : FAIR}
\label{chap:fair}
\chaptermark{Quantifier le risque}

\epigraph{%
    « Un conseil ne vote pas une couleur.\\
    Il vote un budget. »%
}{Principe du cours}

\niveauChapitre{Intermédiaire}{%
    Maîtrise du Chapitre~\ref{chap:iso27005}%
}

\begin{boxobjectifs}
\begin{itemize}
    \item Expliquer pourquoi et quand quantifier un risque cyber.
    \item Décomposer un risque selon la taxonomie FAIR~: fréquence des
          événements de perte et ampleur des pertes.
    \item Produire des estimations calibrées (minimum, plus probable,
          maximum) et en coter la base de preuve.
    \item Conduire une simulation de Monte-Carlo pas à pas.
    \item Interpréter une distribution de pertes annuelles~: moyenne,
          quantiles, probabilité de dépassement.
    \item Comparer des traitements en francs CFA, y compris une assurance.
    \item Dire ce que la quantification apporte et ce qu'elle masque.
\end{itemize}
\end{boxobjectifs}

\bigskip

% ============================================================================
\section*{Avant de commencer : 193 millions, est-ce trop~?}
% ============================================================================

Douala, janvier 2027. Le comité des risques du conseil de la BANQUE
SANAGA\footnote{Comme dans tout le cours, la BANQUE SANAGA, ses
dirigeants, les dates et les montants sont fictifs~; la méthode et les
normes citées sont exactes.} examine le plan de traitement de R-AG-01~:
huit actions, 193~millions FCFA (Chapitre~\ref{chap:iso27005}). Le
directeur financier pose la question que les échelles ne savent pas
traiter~: « Vous me dites que le risque est coté 12. Combien nous
coûte-t-il~? Et combien nous coûtera-t-il après ces 193~millions~? »

La RSSI n'a pas de réponse en francs CFA. La cote~12 classe le risque
parmi les intolérables~; elle ne dit pas si 193~millions sont une bonne
affaire, une dépense insuffisante ou un luxe. Le comité lui demande de
revenir en mars avec une quantification.

\medskip

Ce chapitre présente la méthode qu'elle emploie, FAIR, et la déroule sur
R-AG-01. Tous les résultats numériques ont été obtenus par une simulation
réelle de 100\,000 années, dont le chapitre explique chaque étape.

\begin{boxalert}
\textbf{Violation typique.} Présenter au conseil une perte « attendue » de
143,5~millions, obtenue en multipliant une fréquence par une perte, comme
si c'était ce qui allait arriver.

\smallskip
\textbf{Règle~:} une quantification se présente comme une distribution~:
la plupart des années ne coûtent rien, quelques-unes coûtent très cher. On
donne au moins la moyenne, la probabilité d'une perte et un quantile élevé,
avec les hypothèses et leur base de preuve.
\end{boxalert}

\separateur

% ============================================================================
\section{Pourquoi Quantifier}
\label{sec:c05-pourquoi}
% ============================================================================

Les échelles qualitatives ont deux mérites~: elles sont rapides, et
chacun les comprend. Le Chapitre~\ref{chap:risque} en a montré les
limites~: niveaux ordinaux, produit qui écrase les écarts, frontières
arbitraires. Ces limites deviennent bloquantes dans trois situations~:
\begin{itemize}
    \item \textbf{comparer un risque à une dépense}~: un plan de
          193~millions, une prime d'assurance, un recrutement~;
    \item \textbf{comparer des risques de nature différente}~: un risque
          cyber et un risque de crédit, que le conseil arbitre avec le même
          budget~;
    \item \textbf{une cote qui chevauche un seuil}~: la fréquence de
          R-AG-01, estimée entre 0,2 et 0,6 par an, traverse les niveaux~2
          et~3 de l'échelle, donc la frontière entre «~à traiter dans
          l'année~» et «~intolérable~» (Chapitre~\ref{chap:risque}).
\end{itemize}

On ne quantifie pas tout. La quantification coûte du temps et des
compétences~; on la réserve aux risques qui pèsent sur une décision
importante, typiquement ceux d'impact critique ou dont la cote traverse un
seuil.

% ============================================================================
\section{La Taxonomie FAIR}
\label{sec:c05-taxonomie}
% ============================================================================

\subsection{Une décomposition}

FAIR (\emph{Factor Analysis of Information Risk}) est un modèle de
quantification du risque d'information. Il est normalisé par The Open
Group sous le nom d'Open FAIR, en deux normes~: la taxonomie du risque
(O-RT, version~3.0.1) et l'analyse du risque (O-RA,
version~2.0.1)~\cite{opengroup2026openfair}. Le cours présente la
taxonomie d'après l'ouvrage de référence de Freund et
Jones~\cite{freund2015fair}.

\begin{boxdef}
\termecle{Risque (au sens de FAIR)}\index{FAIR}~: fréquence et ampleur
probables des pertes futures. Il se décompose en deux branches~:
\begin{itemize}
    \item la \textbf{fréquence des événements de perte}~: combien de fois
          par an un événement de menace aboutit à une perte~;
    \item l'\textbf{ampleur des pertes}~: combien coûte chaque événement.
\end{itemize}
\end{boxdef}

\begin{figure}[H]
\centering
\begin{tikzpicture}[
    noeud/.style={draw=CouleurPrimaire, thick, rounded corners,
                  fill=CouleurPrimaire!8, text width=2.3cm,
                  minimum height=0.9cm, align=center, font=\scriptsize},
    racine/.style={noeud, fill=CouleurAccent!15, draw=CouleurAccent},
    lien/.style={thick, CouleurSecondaire}]
    \node[racine] (R) at (0,0) {\textbf{Risque}};
    \node[noeud] (F) at (-3.4,-1.5) {Fréquence des événements de perte};
    \node[noeud] (A) at (3.4,-1.5) {Ampleur des pertes};
    \node[noeud] (T) at (-5.2,-3.1) {Fréquence des événements de menace};
    \node[noeud] (V) at (-1.8,-3.1) {Vulnérabilité};
    \node[noeud] (C) at (-6.0,-4.7) {Fréquence de contact};
    \node[noeud] (P) at (-3.4,-4.7) {Probabilité d'action};
    \node[noeud] (Cap) at (-1.0,-4.7) {Capacité de la menace};
    \node[noeud] (Res) at (1.4,-4.7) {Force de résistance};
    \node[noeud] (PP) at (2.0,-3.1) {Pertes primaires};
    \node[noeud] (PS) at (4.8,-3.1) {Pertes secondaires};
    \draw[lien] (R) -- (F); \draw[lien] (R) -- (A);
    \draw[lien] (F) -- (T); \draw[lien] (F) -- (V);
    \draw[lien] (T) -- (C); \draw[lien] (T) -- (P);
    \draw[lien] (V) -- (Cap); \draw[lien] (V) -- (Res);
    \draw[lien] (A) -- (PP); \draw[lien] (A) -- (PS);
\end{tikzpicture}
\caption{La taxonomie FAIR (d'après Freund et Jones, 2015)}
\label{fig:c05-taxonomie}
\end{figure}

\begin{tableaucours}{L{3.6cm}L{8.6cm}}{%
    \tableauHeader{Facteur} &
    \tableauHeader{Définition}
}
Fréquence des événements de menace & Nombre de fois par an où une source
de menace agit contre l'actif~: produit de la fréquence de contact et de
la probabilité qu'elle passe à l'action \\
Vulnérabilité & Probabilité qu'un événement de menace devienne un
événement de perte~: elle dépend de la capacité de la menace comparée à la
force de résistance des mesures \\
Pertes primaires & Pertes subies directement par l'organisation du fait de
l'événement \\
Pertes secondaires & Pertes qui viennent des réactions de tiers~: clients,
régulateur, tribunaux, presse~; elles ont leur propre fréquence et leur
propre ampleur \\
\end{tableaucours}
\vspace{-0.8em}
\begin{center}{\small\itshape Tableau~5.1. Les facteurs de la taxonomie}\end{center}

\begin{boxinfo}
\textbf{À vérifier dans la norme.} La version~3 de la taxonomie Open FAIR
a révisé certains termes~; on y lit notamment \emph{susceptibilité} là où
les textes antérieurs parlent de \emph{vulnérabilité}. Le cours garde le
vocabulaire de Freund et Jones, plus répandu, et signale l'écart.
\end{boxinfo}

\subsection{Six formes de pertes}

FAIR ventile l'ampleur en six formes de pertes~: la \emph{productivité}
(activité perdue), la \emph{réponse} (gestion de l'incident,
communication, conseils), le \emph{remplacement} (actifs à reconstruire
ou racheter), les \emph{amendes et jugements}, l'\emph{avantage
concurrentiel} perdu, et la \emph{réputation}~\cite{freund2015fair}.

\begin{boxscenario}{Opération SANAGA. Les 410 millions, ventilés}
\textbf{Constatation.} La direction financière ventile le coût de
l'incident de septembre 2026~:
\begin{itemize}
    \item productivité~: 190~millions (neuf agences fermées trois jours,
          opérations perdues ou reportées)~;
    \item réponse~: 140~millions (prestataire de réponse à incident,
          cellule de crise, communication auprès de la clientèle)~;
    \item remplacement~: 80~millions (reconstruction de 180 postes)~;
    \item amendes et jugements, avantage concurrentiel~: néant~;
    \item réputation~: non chiffrée.
\end{itemize}

\textbf{Analyse} (confiance élevée sur les pertes primaires, faible sur
les secondaires). Les 410~millions sont presque entièrement des pertes
primaires. La perte de clients qui ont changé de banque n'est pas mesurée~:
elle est \NEK{}. Le coût réel de l'incident est donc d'au moins
410~millions, pas de 410~millions.
\end{boxscenario}

% ============================================================================
\section{Estimer}
\label{sec:c05-estimer}
% ============================================================================

\subsection{Des estimations calibrées}

FAIR ne demande pas des valeurs exactes, que personne n'a, mais des
\termecle{estimations calibrées}\index{estimation calibrée}~: pour chaque
facteur, un minimum, une valeur la plus probable et un maximum, choisis de
sorte que l'expert soit confiant, à 90~\% environ, que la vraie valeur se
trouve entre le minimum et le maximum. Hubbard et Seiersen montrent que
cette confiance s'apprend, par des exercices de calibrage, et que la
plupart des experts non entraînés donnent des intervalles trop
étroits~\cite{hubbard2016measure}.

\subsection{Les estimations de R-AG-01}

En février 2027, la RSSI anime un atelier avec la direction financière, la
direction du réseau d'agences et le prestataire de réponse à incident.

\begin{tableaucours}{L{3.0cm}ccc L{4.4cm}}{%
    \tableauHeader{Facteur} &
    \tableauHeader{Min.} &
    \tableauHeader{Plus prob.} &
    \tableauHeader{Max.} &
    \tableauHeader{Justification}
}
Fréquence des événements de perte (par an) & 0,2 & 0,35 & 0,6 & Un
incident réussi en 2026~; tentatives fréquentes dans le secteur~; défenses
de 2025 \\
Ampleur par événement (M FCFA) & 150 & 410 & 900 & Minimum~: trois agences
touchées~; plus probable~: l'incident vécu~; maximum~: propagation à tout
le réseau et pertes secondaires \\
\end{tableaucours}
\vspace{-0.8em}
\begin{center}{\small\itshape Tableau~5.2. Estimations calibrées pour
R-AG-01 (défenses de 2025)}\end{center}

Pour la fréquence, l'atelier estime directement la fréquence des
événements de perte, sans descendre jusqu'à la fréquence de contact et à la
vulnérabilité~: la taxonomie permet de s'arrêter au niveau où l'on dispose
de la meilleure information.

\subsection{La base de preuve des estimations}

\begin{boxlecture}{Base de preuve des six valeurs}
\textbf{Constatation.} La RSSI cote la base de chacune des six valeurs
du Tableau~5.2~:
\begin{itemize}
    \item fréquence~: minimum \NEK{} (avis non motivé), plus probable
          \OUI{} (incident de 2026, données sectorielles d'un assureur),
          maximum \NEK{}~;
    \item ampleur~: minimum \OUI{} (devis de reconstruction pour trois
          agences), plus probable \OUI{} (coût de l'incident), maximum
          \NEK{} (pertes secondaires non chiffrées).
\end{itemize}
3~\OUI{}, 0~\NON{}, 3~\NEK{}~: $\mu = 0{,}50$~; $U = 0{,}50$~;
$H = -(2 \times 0{,}5\log_2 0{,}5) = 1{,}00$~bit.

\medskip
\textbf{Analyse} (confiance élevée). La quantification est solide en son
centre et incertaine dans ses queues. Or ce sont les queues qui
intéressent le conseil~: le maximum d'ampleur décide du quantile à
99~\%, et le maximum de fréquence de la probabilité d'avoir deux incidents
la même année. La Section~\ref{sec:c05-masque} mesure ce que pèse cette
incertitude.

\medskip
\hfill\lecture{0{,}50}{0{,}50}{1{,}00}
\end{boxlecture}

% ============================================================================
\section{La Simulation de Monte-Carlo}
\label{sec:c05-montecarlo}
% ============================================================================

\subsection{Le principe}

Multiplier la fréquence la plus probable par l'ampleur la plus probable
donne $0{,}35 \times 410 = 143{,}5$~millions. Ce chiffre n'est ni ce qui
arrivera une année donnée (le plus souvent, rien), ni même la moyenne,
parce qu'il ignore la forme des distributions. La \termecle{simulation de
Monte-Carlo}\index{simulation de Monte-Carlo} remplace ce calcul par un
grand nombre d'années tirées au hasard, selon les estimations.

\begin{boxprocedure}
\textbf{Simuler une année.}
\begin{enumerate}
    \item Tirer une fréquence $\lambda$ dans la distribution de fréquence
          (ici une loi PERT, qui suit le minimum, le plus probable et le
          maximum).
    \item Tirer un nombre d'événements dans l'année, selon une loi de
          Poisson de paramètre $\lambda$~: le plus souvent 0, parfois 1,
          rarement 2 ou plus.
    \item Pour chaque événement, tirer une ampleur dans la distribution
          d'ampleur (loi PERT).
    \item La perte de l'année est la somme des ampleurs.
\end{enumerate}
On répète l'opération un grand nombre de fois (ici 100\,000), puis on lit
la distribution des pertes annuelles obtenues. Le papier de travail
conserve le programme, la graine du générateur aléatoire et les
estimations, pour que le résultat puisse être reproduit.
\end{boxprocedure}

\begin{boxinfo}
\textbf{La loi PERT.} Elle se construit à partir d'un minimum $a$, d'une
valeur la plus probable $m$ et d'un maximum $b$. Sa moyenne vaut
$\frac{a + 4m + b}{6}$~: pour la fréquence,
$\frac{0{,}2 + 1{,}4 + 0{,}6}{6} \approx 0{,}367$~; pour l'ampleur,
$\frac{150 + 1\,640 + 900}{6} \approx 448$~millions. La moyenne de
l'ampleur dépasse la valeur la plus probable, parce que le maximum est plus
éloigné du mode que le minimum.
\end{boxinfo}

\subsection{Dix années, à la main}

Pour comprendre la mécanique, voici les dix premières années d'une
simulation.

\begin{tableaucours}{cccL{3.4cm}r}{%
    \tableauHeader{Année} &
    \tableauHeader{$\lambda$ tirée} &
    \tableauHeader{Événements} &
    \tableauHeader{Ampleurs (M FCFA)} &
    \tableauHeader{Perte}
}
1 & 0,38 & 0 & -- & 0 \\
2 & 0,26 & 0 & -- & 0 \\
3 & 0,39 & 0 & -- & 0 \\
4 & 0,43 & 1 & 385 & 385 \\
5 & 0,38 & 0 & -- & 0 \\
6 & 0,26 & 0 & -- & 0 \\
7 & 0,43 & 0 & -- & 0 \\
8 & 0,35 & 0 & -- & 0 \\
9 & 0,40 & 1 & 521 & 521 \\
10 & 0,35 & 0 & -- & 0 \\
\textbf{Moyenne} & & & & \textbf{90,6} \\
\end{tableaucours}
\vspace{-0.8em}
\begin{center}{\small\itshape Tableau~5.3. Dix années simulées}\end{center}

Huit années sur dix ne coûtent rien~; deux coûtent 385 et 521~millions.
La moyenne de ces dix années, 90,6~millions, est loin de la moyenne
obtenue sur 100\,000 années (164,6~millions)~: dix tirages ne suffisent
pas. C'est la leçon de la simulation, et aussi celle de l'expérience~: une
banque qui n'a pas eu d'incident pendant huit ans n'en déduit rien sur la
neuvième.

\subsection{Cent mille années}

\begin{tableaucours}{L{6.6cm}rr}{%
    \tableauHeader{Indicateur (M FCFA, par an)} &
    \tableauHeader{Défenses 2025} &
    \tableauHeader{Après traitement}
}
Perte annuelle moyenne & 164,6 & 18,6 \\
Probabilité d'une perte dans l'année & 30,5~\% & 10,1~\% \\
Médiane (une année sur deux) & 0 & 0 \\
Quantile à 90~\% (une année sur dix au-delà) & 579 & 60 \\
Quantile à 95~\% (une année sur vingt) & 734 & 174 \\
Quantile à 99~\% (une année sur cent) & 1\,170 & 283 \\
Probabilité d'une perte supérieure à 500~M & 14,3~\% & 0,06~\% \\
\end{tableaucours}
\vspace{-0.8em}
\begin{center}{\small\itshape Tableau~5.4. Distribution des pertes annuelles
dues à R-AG-01 (100\,000 années simulées)}\end{center}

La colonne «~après traitement~» reprend la simulation avec les estimations
de l'atelier pour la situation visée par le plan~: fréquence entre 0,05 et
0,20 (plus probable 0,10), grâce au second facteur et au cloisonnement~;
ampleur entre 50 et 400~millions (plus probable 150), grâce aux
sauvegardes testées et au mode dégradé.

\begin{tableaucours}{cccccc}{%
    \tableauHeader{Perte supérieure à (M)} &
    \tableauHeader{0} &
    \tableauHeader{250} &
    \tableauHeader{500} &
    \tableauHeader{750} &
    \tableauHeader{1\,000}
}
Défenses 2025 & 30,5~\% & 28,8~\% & 14,3~\% & 4,7~\% & 2,1~\% \\
Après traitement & 10,1~\% & 1,8~\% & 0,06~\% & 0~\% & 0~\% \\
\end{tableaucours}
\vspace{-0.8em}
\begin{center}{\small\itshape Tableau~5.5. Probabilités de dépassement
(courbe de dépassement, en points)}\end{center}

% ============================================================================
\section{Interpréter et Décider}
\label{sec:c05-decider}
% ============================================================================

\subsection{Lire la distribution}

\begin{boxlecture}{Ce que la RSSI présente au comité, en mars 2027}
\textbf{Constatation.}
\begin{itemize}
    \item Avec les défenses de 2025, une année sur trois environ connaît
          un incident~; la perte moyenne est de 165~millions par an, et une
          année sur sept dépasse 500~millions.
    \item Après traitement, une année sur dix connaît un incident~; la
          perte moyenne tombe à 19~millions, et une perte de plus de
          500~millions devient exceptionnelle (moins d'une année sur mille).
\end{itemize}

\textbf{Analyse} (confiance moyenne~: voir la base de preuve des
estimations). Le plan réduit la perte moyenne d'environ 146~millions par
an. Il coûte 193~millions la première année, puis environ 30~millions par
an (licences, tests trimestriels, exercices). Le gain net annuel, de
l'ordre de 116~millions, rembourse l'investissement en moins de deux ans.
Sur ce seul risque, 193~millions ne sont pas un luxe.
\end{boxlecture}

\subsection{L'appétence en francs CFA}

La quantification permet aussi d'exprimer l'appétence autrement qu'en
couleurs. La RSSI propose au comité une formulation~: « la probabilité
qu'un même risque cyber cause, une année donnée, une perte supérieure à
500~millions FCFA ne dépasse pas 5~\% ». Avec les défenses de 2025,
R-AG-01 la viole largement (14,3~\%)~; après traitement, il la respecte
(0,06~\%).

La simulation tranche aussi la question laissée ouverte au
Chapitre~\ref{chap:risque}~: la fréquence tirée dépasse 1/3 par an dans
64~\% des années simulées. Le niveau~3 de vraisemblance est donc plus
probable que le niveau~2~; la cote prudente de 12 était justifiée.

\subsection{L'assurance, chiffrée}

L'offre du courtier (prime de 35~millions par an, franchise de
100~millions par événement, plafond de 1~milliard) se juge maintenant sur
les mêmes simulations.

\begin{tableaucours}{L{6.4cm}rr}{%
    \tableauHeader{M FCFA par an} &
    \tableauHeader{Défenses 2025} &
    \tableauHeader{Après traitement}
}
Perte moyenne sans assurance & 164,6 & 18,6 \\
Remboursement moyen de l'assureur & 126,6 & 8,2 \\
Coût moyen avec assurance (perte $-$ remboursement $+$ prime) & 73,0 &
45,4 \\
\end{tableaucours}
\vspace{-0.8em}
\begin{center}{\small\itshape Tableau~5.6. L'assurance, avant et après
traitement}\end{center}

Avant traitement, l'assurance divise le coût moyen par plus de deux. Après
traitement, elle coûte en moyenne plus qu'elle ne rapporte (45,4 contre
18,6). Deux conclusions~: l'assurance est utile comme protection
temporaire, pendant la mise en œuvre du plan~; ensuite, elle ne se justifie
que si la prime baisse, ou pour protéger la banque contre les années
extrêmes que la simulation estime mal. Et l'on ne pourra pas l'obtenir sans
les mesures, que les assureurs exigent~: la réduction vient d'abord.

% ============================================================================
\section{Ce que la Quantification Masque}
\label{sec:c05-masque}
% ============================================================================

\subsection{Des chiffres faits d'opinions}

Une simulation transforme des estimations en distribution~; elle ne les
rend pas plus vraies. Trois des six valeurs d'entrée sont \NEK{}. Pour
mesurer ce que cela pèse, la RSSI relance la simulation en portant le
maximum d'ampleur de 900 à 1\,500~millions, pour tenir compte des pertes
secondaires non chiffrées.

\begin{tableaucours}{L{6.6cm}rr}{%
    \tableauHeader{Indicateur (M FCFA)} &
    \tableauHeader{Max. 900} &
    \tableauHeader{Max. 1\,500}
}
Perte annuelle moyenne & 164,6 & 200,9 \\
Quantile à 99~\% & 1\,170 & 1\,507 \\
Probabilité d'une perte supérieure à 500~M & 14,3~\% & 18,4~\% \\
\end{tableaucours}
\vspace{-0.8em}
\begin{center}{\small\itshape Tableau~5.7. Sensibilité au maximum
d'ampleur}\end{center}

Une seule valeur \NEK{}, modifiée, déplace la moyenne de 22~\% et le
quantile à 99~\% de 29~\%. La décision de traiter ne change pas~; mais le
montant de couverture d'assurance, lui, en dépend. C'est le rôle de
l'analyse de sensibilité~: identifier les \NEK{} dont la décision dépend,
et y consacrer l'effort de collecte de preuves.

\subsection{Cinq angles morts}

\begin{itemize}
    \item \textbf{Le scénario absent.} On ne simule que ce que l'on a
          imaginé. En 2025, R-AG-01 n'existait pas dans le registre~:
          aucune simulation ne l'aurait chiffré.
    \item \textbf{La fausse précision.} « 164,6~millions » suggère une
          exactitude que les estimations n'ont pas. On présente des ordres
          de grandeur arrondis («~environ 165~millions~»).
    \item \textbf{L'indépendance supposée.} Le modèle tire les événements
          indépendamment. Une même faiblesse peut provoquer plusieurs
          incidents corrélés, et une crise du secteur toucher plusieurs
          banques à la fois.
    \item \textbf{L'adversaire qui s'adapte.} Les estimations décrivent la
          menace de 2027~; un attaquant qui change de méthode les rend
          caduques (Chapitre~\ref{chap:risque}, Tableau~2.2).
    \item \textbf{L'incertitude de connaissance.} La distribution mêle
          l'aléa (l'attaque aura-t-elle lieu~?) et l'ignorance (la vraie
          fréquence)~; elle ne dit pas quelle part vient de chacun. La cote
          de la base de preuve des estimations le dit.
\end{itemize}

\begin{boxinfo}
\textbf{La quantification et l'auditologie.} Le cours tient les deux
ensemble~: la simulation décrit l'incertitude \emph{aléatoire} sur les
pertes, sous des hypothèses~; la cote \OUI{}~/ \NON{}~/ \NEK{} des
hypothèses décrit l'incertitude \emph{épistémique} sur ces hypothèses. Une
quantification présentée sans la seconde est incomplète~; le
Chapitre~\ref{chap:auditologie} montre comment les combiner.
\end{boxinfo}

\begin{boxinfo}
\textbf{Et le cadre juridique~?} Aucun texte camerounais n'impose de
quantifier les risques cyber. L'article~26 de la loi n°~2010/012 demande
de les \emph{évaluer} (Chapitre~\ref{chap:iso27005})~\cite{loi2010012}~;
la quantification est une manière de le faire, que le conseil, le
superviseur et l'assureur comprennent, parce qu'elle parle leur langue~:
l'argent.
\end{boxinfo}

% ============================================================================
\section{Ce que la Quantification demande à chaque Métier}
\label{sec:c05-metiers}
% ============================================================================

\begin{boxmetier}{Gestionnaire des risques cyber}
\textbf{Quantifier ce qui décide}
\begin{enumerate}
    \item Réserver la quantification aux risques dont dépend une décision
          importante.
    \item Faire produire des estimations calibrées, par des experts
          entraînés, et coter leur base de preuve.
    \item Présenter une distribution (moyenne, probabilité, quantile élevé),
          jamais un seul chiffre.
    \item Tester la sensibilité aux valeurs \NEK{}.
\end{enumerate}
\end{boxmetier}

\begin{boxmetier}{RSSI}
\textbf{Défendre un budget}

La RSSI ne demande plus 193~millions pour « réduire un risque coté 12 ».
Elle montre une réduction de perte moyenne de l'ordre de 146~millions par
an, et le passage de la probabilité d'une année à plus de 500~millions de
14~\% à moins de 0,1~\%. Le conseil vote un budget, avec ses hypothèses.
\end{boxmetier}

\begin{boxmetier}{Auditeur en cybersécurité}
\textbf{Auditer les entrées, pas le logiciel}

L'auditeur ne refait pas les 100\,000 tirages~; il vérifie que le programme
et la graine sont conservés, que chaque estimation a une justification
écrite, que les auteurs ont été calibrés, et que la sensibilité aux
valeurs non fondées a été testée.
\end{boxmetier}

\begin{boxmetier}{Responsable GRC}
\textbf{Une même monnaie}

Le responsable GRC propose au conseil une appétence exprimée en francs
CFA, applicable à tous les risques quantifiés, cyber ou non. Les risques
deviennent comparables, et les arbitrages budgétaires, explicables.
\end{boxmetier}

\separateur

\begin{boxresume}
\textbf{Ce qu'il faut retenir de ce chapitre~:}
\begin{itemize}
    \item \textbf{Quantifier} sert à comparer un risque à une dépense, des
          risques entre eux, ou à trancher une cote qui chevauche un seuil.
    \item \textbf{FAIR} décompose le risque en fréquence des événements de
          perte et ampleur des pertes~; l'ampleur se ventile en six formes
          de pertes, primaires et secondaires.
    \item \textbf{Des estimations calibrées} (minimum, plus probable,
          maximum) remplacent les valeurs exactes que personne n'a~; leur
          base de preuve se cote.
    \item \textbf{La simulation de Monte-Carlo} tire des années~; on lit la
          moyenne, la probabilité d'une perte et les quantiles élevés.
    \item \textbf{R-AG-01}~: environ 165~millions de perte moyenne par an
          avant traitement, 19 après~; le plan de 193~millions se rembourse
          en moins de deux ans.
    \item \textbf{La quantification masque} le scénario absent, la fausse
          précision, les corrélations, l'adversaire qui s'adapte, et
          l'incertitude de connaissance.
\end{itemize}
\end{boxresume}

\begin{boxinfo}
\textbf{Points de contrôle.} Vous avez compris ce chapitre si vous savez~:
\begin{itemize}
    \item placer un facteur dans la taxonomie FAIR~;
    \item ventiler un coût d'incident en formes de pertes~;
    \item simuler quelques années à la main et dire pourquoi elles ne
          suffisent pas~;
    \item lire une distribution et en tirer une décision de traitement ou
          d'assurance.
\end{itemize}
\end{boxinfo}

% ============================================================================
\section*{Exercices}
% ============================================================================

\begin{boxexo}[title={\ding{45}\quad Exercice~5.1. Six formes de pertes \niveaubase}]
Classez chacun des coûts suivants dans l'une des six formes de pertes de
FAIR, et dites s'il s'agit d'une perte primaire ou secondaire.
\begin{enumerate}[(a)]
    \item Les heures supplémentaires des agents pour ressaisir les
          opérations.
    \item Le rachat de vingt postes de travail.
    \item L'amende infligée par le superviseur pour notification tardive.
    \item Les honoraires du prestataire de réponse à incident.
    \item Les clients qui ferment leur compte après l'incident.
    \item Le contrat d'une grande entreprise remporté par une banque
          concurrente pendant la fermeture des guichets.
\end{enumerate}
\end{boxexo}

\begin{boxexo}[title={\ding{45}\quad Exercice~5.2. Le chiffre trompeur \niveaumoyen}]
\begin{enumerate}[(a)]
    \item Calculez les moyennes PERT de la fréquence et de l'ampleur pour la
          situation après traitement (Section~\ref{sec:c05-montecarlo}), et
          leur produit. Comparez au résultat de la simulation (18,6).
    \item Pourquoi la médiane de la perte annuelle est-elle nulle, avant
          comme après traitement~?
    \item Un administrateur propose de provisionner chaque année
          165~millions «~pour couvrir le risque~». Expliquez ce que cette
          provision couvre et ce qu'elle ne couvre pas, à l'aide du
          Tableau~5.4.
\end{enumerate}
\end{boxexo}

\begin{boxexo}[title={\ding{45}\quad Exercice~5.3. Une autre offre d'assurance \niveaumoyen}]
Un second assureur propose, après traitement, une prime de 12~millions par
an, une franchise de 200~millions par événement et un plafond de
500~millions. On admet qu'après traitement il rembourserait en moyenne
1,5~million par an.
\begin{enumerate}[(a)]
    \item Calculez le coût moyen annuel avec cette assurance et comparez-le
          à l'absence d'assurance et à la première offre.
    \item À l'aide du Tableau~5.5, dites dans quelles années cette
          assurance joue, et ce qu'elle protège réellement.
    \item Recommandez une décision au comité, en deux phrases.
\end{enumerate}
\end{boxexo}

\begin{boxexo}[title={\ding{45}\quad Exercice~5.4. Quantifier SO1 pour PHARMACAM \niveauexpert}]
LiZbiZ veut présenter à PHARMACAM une quantification simple du scénario
SO1 (Chapitre~\ref{chap:ebios}). L'atelier retient~: un événement au plus
par an, avec une probabilité de 0,25~; une ampleur triangulaire de minimum
100, plus probable 300, maximum 800~millions FCFA (moyenne~: le tiers de la
somme des trois valeurs).
\begin{enumerate}[(a)]
    \item Calculez la perte annuelle moyenne.
    \item Simulez dix années à la main avec les nombres aléatoires suivants
          (un événement survient si le nombre est inférieur à 0,25)~: 0,81~;
          0,12~; 0,55~; 0,37~; 0,09~; 0,66~; 0,93~; 0,21~; 0,48~; 0,70.
          Les ampleurs tirées pour les événements sont, dans l'ordre, 320,
          610 et 180~millions. Calculez la moyenne des dix années et
          commentez.
    \item Cotez la base de preuve des quatre valeurs d'entrée, sachant que
          LiZbiZ n'a jamais subi de rançongiciel et que l'ampleur
          maximale correspond à l'estimation de PHARMACAM pour un camion de
          vaccins perdu.
    \item Rédigez le paragraphe destiné à PHARMACAM~: ce que la
          quantification dit, ce qu'elle suppose, ce qu'elle ne dit pas.
\end{enumerate}
\end{boxexo}

\bigskip

\begin{boxinfo}
\textbf{Pour aller plus loin.}
\begin{itemize}
    \item Freund et Jones, \emph{Measuring and Managing Information
          Risk}~\cite{freund2015fair}, l'exposé de référence de FAIR.
    \item Les normes Open FAIR de The Open Group~\cite{opengroup2026openfair}.
    \item Hubbard et Seiersen, \emph{How to Measure Anything in
          Cybersecurity Risk}~\cite{hubbard2016measure}, sur le calibrage
          des experts.
\end{itemize}
\end{boxinfo}

\begin{boxinfo}
\textbf{Transition.} La Partie~II a jusqu'ici appris à apprécier, traiter
et chiffrer les risques. Reste à vérifier que les mesures annoncées
existent et fonctionnent~: c'est le travail de l'auditeur. Le
Chapitre~\ref{chap:audit} présente la conduite d'un audit de
cybersécurité, du mandat au rapport.
\end{boxinfo}
